\documentclass{article}
\usepackage{amsmath, amssymb, amsthm}
\usepackage{hyperref}
\usepackage{geometry}
\geometry{margin=1in}

\title{Erd\H{o}s Problem \#665}
\date{Last updated: 2025-08-31}
\author{Source: erdosproblems.com}

\begin{document}
\maketitle

\noindent\textbf{Status:} OPEN \quad \textbf{No prize}

\noindent\textbf{Tags:} combinatorics

\medskip

\noindent\textbf{Problem Statement:}

A pairwise balanced design for $\{1,\ldots,n\}$ is a collection of sets $A_1,\ldots,A_m\subseteq \{1,\ldots,n\}$ such that $2\leq \lvert A_i\rvert &lt;n$ and every pair of distinct elements $x,y\in \{1,\ldots,n\}$ is contained in exactly one $A_i$. Is there a constant $C&gt;0$ and, for all large $n$, a pairwise balanced design such that\[\lvert A_i\rvert &#62; n^{1/2}-C\]for all $1\leq i\leq m$?




A problem of Erd\H{o}s and Larson \cite{ErLa82}. In general, as Erd\H{o}s asks in \cite{Er97f}, find the slowest growing function $h$ such that, for all large $n$, there exists a pairwise balanced design with\[\lvert A_i\rvert &gt; n^{1/2}-h(n)\]for all $1\leq i\leq m$. The problem above asks whether $h(n)\ll 1$. Erd\H{o}s and Larson prove that $h(n) \ll n^{1/2-c}$ for some constant $c&gt;0$, and note this can be improved to $h(n)\ll (\log n)^2$ assuming Cramer-type bounds on the difference between consecutive primes. Shrikhande and Singhi \cite{ShSi85} have proved that the answer is no conditional on the conjecture that the order of every projective plane is a prime power (see [723] ), by proving that every pairwise balanced design on $n$ points in which each block is of size $\geq n^{1/2}-c$ can be embedded in a projective plane of order $n+i$ for some $i\leq c+2$, if $n$ is sufficiently large. In general, if $H(n)$ is the largest prime gap $\leq n$, then the above reuslts show that, assuming the prime power conjecture, $H(n)\asymp h(n)$.




References

[Er97f] Erd\H{o}s, Paul, Some unsolved problems . Combinatorics, geometry and probability (Cambridge, 1993) (1997), 1-10.

[ErLa82] Erd\H{o}s, P. and Larson, J., On pairwise balanced block designs with the sizes of blocks as uniform as possible . Annals of Discrete Mathematics (1982), 129-134.

[ShSi85] S. S. Shrikhande and N. M. Singhi, On a problem of Erd\H{o}s and Larson . Combinatorica (1985), 351-358.


\bigskip
\noindent\small{Source: \url{https://www.erdosproblems.com/665}}
\end{document}
